\chapter{Related Work}

This thesis sits at the boundary between replicated transaction systems and
consensus engineering. It asks whether an existing high-throughput transaction
system can use Raft as its replication backend without losing the performance
shape that made the original system fast. The closest related work therefore
falls into five groups: Mako and ROLIS, replicated transaction systems,
consensus protocols, leader placement and replay parallelism, and
storage-aware evaluation.

\section{Mako, ROLIS, and Execute-Replicate-Replay}

Mako is the system extended in this thesis \cite{mako-osdi25}. It is a
geo-replicated transactional key-value store whose original replication
backend is Paxos-based. Its performance depends on the
execute-replicate-replay pipeline used throughout this thesis: leader-side
workers execute transactions and produce log records, the replication layer
orders those records, and follower replicas replay committed records into
their local state.

ROLIS provides the closest conceptual predecessor for this pipeline
\cite{rolis-eurosys22}. Its central lesson is that high-throughput replicated
transactions cannot treat replication as a small final step after execution.
Replication, replay, and multicore execution must be engineered together. This
thesis applies that lesson to a concrete backend substitution inside Mako:
replacing the original Paxos path with Raft while preserving the surrounding
pipeline.

That makes Mako and ROLIS different from generic consensus references in this
thesis. They define the performance problem. The relevant question is not
only whether replicas agree on a log, but whether agreement can happen without
collapsing the multicore execution and replay structure above it.

\section{Replicated and Geo-Distributed Transactions}

Several distributed transaction systems also study how to combine strong
consistency, replication, and high throughput. Spanner is a landmark
geo-distributed database that uses synchronous replication and TrueTime to
provide externally consistent distributed transactions at global scale
\cite{spanner-osdi12}. Spanner's design center is a complete global database
architecture with time-based consistency and production-scale operation; this
thesis focuses on a narrower backend replacement inside Mako.

Calvin uses deterministic transaction ordering to reduce distributed
coordination costs in partitioned databases \cite{calvin-sigmod12}. It is
relevant because it also treats ordering as part of transaction processing,
but it orders transaction inputs ahead of execution. Mako instead executes
speculatively and replicates transaction log records afterward.

TAPIR and Janus revisit the boundary between transaction processing and
replication. TAPIR builds a transaction protocol over inconsistent
replication, moving the consistency burden into the transaction layer
\cite{tapir-sosp15}. Janus consolidates concurrency control and consensus for
contended wide-area transactions \cite{janus-osdi16}. These systems show that
the transaction/replication boundary is often where important performance
decisions live.

This thesis is more conservative than those systems. It does not replace
Mako's transaction protocol or redesign concurrency control. Instead, it keeps
the transaction layer fixed and asks how far a standard leader-based consensus
protocol can go when the surrounding integration is engineered carefully.

\section{Consensus Protocols and Consensus Engineering}

Paxos and Raft are the main consensus protocols behind this thesis
\cite{paxos-made-simple,raft-atc14}. Paxos is the original Mako backend. Raft
is attractive because it packages leader election, log replication, and leader
changes in a protocol designed to be easier to understand and implement. This
thesis compares them as systems backends rather than as abstract protocols.

Paxos Made Live is especially relevant because it argues, from engineering
experience, that turning a consensus algorithm into a working replicated
system requires many practical decisions beyond the pseudocode
\cite{paxos-made-live}. This thesis is in the same spirit: the integration
must decide where leaders should be placed, how commands should be routed, how
followers should replay, and how persistence costs should be measured.

Other consensus protocols, such as EPaxos, improve the protocol structure
itself by reducing reliance on a single stable leader and balancing load
across replicas \cite{epaxos-sosp13}. This thesis takes a different path: it
keeps ordinary Raft as the consensus protocol and adapts the surrounding Mako
integration through topology, leader placement, and replay offload.

This contrast is useful because it separates two ways to address bottlenecks.
One approach is to change the consensus protocol. The approach here is to keep
Raft understandable and change the system around it so that Raft fits the
application pipeline.

\section{Leader Placement and Replay Parallelism}

Leader placement and replay parallelism are systems issues that consensus
alone does not solve. A leader-based protocol may be safe no matter which
eligible replica leads, but application performance can still depend on where
that leader is located. Similarly, a consensus protocol can decide an order
for commands while leaving the application responsible for applying committed
work efficiently.

ROLIS and Mako both show that replay must be treated as part of the
performance pipeline \cite{rolis-eurosys22,mako-osdi25}. This thesis extends
that lesson to Single-Raft: after many worker lanes feed one consolidated Raft
log, follower replay has to be parallelized outside the consensus protocol.

The preferred-leader mechanism fits the same pattern. Raft safety does not
require a particular eligible leader, but Mako's submission path performs best
when leadership is placed where transaction records are produced. The
mechanism is therefore best understood as application-aware leader placement
around Raft, not as a new consensus rule.

\section{Persistence and Storage-Aware Evaluation}

Durability is also an engineering concern rather than only a protocol detail.
Replicated systems must persist enough state to recover safely, and the cost
of that persistence can dominate throughput. Paxos Made Live highlights the
practical complexity of durable consensus state \cite{paxos-made-live}, and
Spanner illustrates that replication, durability, and production storage
behavior are part of a complete database system \cite{spanner-osdi12}.

This thesis uses controlled simulated-disk settings and records bytes, writes,
and normalized storage work charged to each backend. The goal is not to model
every storage-stack detail, but to make persistence results explainable rather
than relying on throughput curves alone.
